\documentclass{article}
\usepackage{tikz,amsmath,siunitx}
\usetikzlibrary{arrows,snakes,backgrounds,patterns,matrix,shapes,fit,calc,shadows,plotmarks}
\usepackage[graphics,tightpage,active]{preview}
\usepackage{xcolor}
\usepackage{ulem}
\usepackage{amssymb}
\newcommand\Scale{3}
\newcommand\HEIGHT{54.5}
\newcommand\WIDTH{90}
\newcommand\HSPACEO{0.15cm}
\newcommand\HSPACET{0.22cm}

\newcommand\INDENT{1.00cm}

\PreviewEnvironment{tikzpicture}
\PreviewEnvironment{equation}
\PreviewEnvironment{equation*}
\newlength{\imagewidth}
\newlength{\imagescale}
\pagestyle{empty}
\thispagestyle{empty}
\begin{document}
%用于二次互反律的证明.
\begin{tikzpicture}[
    thick,
    >=stealth',
    dot/.style = {
      draw,
      fill = white,
      circle,
      inner sep = 0pt,
      minimum size = 4pt
    }
  ]

%坐标系
\coordinate (O) at (0,0);
  \draw[->] (-0.3,0) -- (8,0) ;
  \draw[->] (0,-0.3) -- (0,5) ;
  \draw (O) -- (7, 3.77);
  \draw(0, 3.5) -- (6.5, 3.5) -- (6.5, 0 );
  \node [below left] at (O) {$(0, 0)$};
  \node [below right] at (6.5, 3.5) {($p/2, q/2$)};
  \node [below] at (6.5, 0) {($p/2, 0$)};
  \node [left] at (0, 3.5) {($0, q/2$)};
  
  \draw[fill] (0,0) circle [radius=0.05];
   \draw[fill] (0.5,0) circle [radius=0.05];
   \draw[fill] (1,0) circle [radius=0.05];
   \draw[fill] (1.5,0) circle [radius=0.05];
   \draw[fill] (2,0) circle [radius=0.05];
   
    \draw[fill] (0,0.5) circle [radius=0.05];
     \draw[fill] (0,1) circle [radius=0.05];
     
     \draw[fill] (0.5,0.5) circle [radius=0.05];
     \draw[fill] (1, 0.5) circle [radius=0.05];
     \draw[fill] (1.5, 0.5) circle [radius=0.05];
      \draw[fill] (2, 0.5) circle [radius=0.05];
      \draw[fill] (2.5, 0.5) circle [radius=0.05];
      \draw[fill] (3, 0.5) circle [radius=0.05];
       \draw[fill] (3.5 , 0.5) circle [radius=0.05];
       
        \draw[fill] (0.5,  1) circle [radius=0.05];
        \draw[fill] (1,  1) circle [radius=0.05];
        \draw[fill] (1.5, 1) circle [radius=0.05];
         \draw[fill] (2, 1) circle [radius=0.05];
         \draw[fill] (2.5, 1) circle [radius=0.05];
         \draw[fill] (3, 1) circle [radius=0.05];
         \draw[fill] (3.5 , 1) circle [radius=0.05];
          
          \draw[fill] (4.5 , 0) circle [radius=0.05];
          \node [below] at (4.5 , 0) {$(k, 0)$};
          \draw[fill] (4.5 , 0.5) circle [radius=0.05];
          \draw[fill] (4.5 , 1) circle [radius=0.05];
          \draw[fill] (4.5 , 1.5) circle [radius=0.05];
          \draw[fill] (4.5 , 2) circle [radius=0.05];
          \draw[fill] (4.5 , 2.5) circle [radius=0.05];
          \node [right] at (4.5, 2) {$(k, \left \lfloor kq/p \right\rfloor)$};
          
          \node  at (3, 2) {$D$};
          \node [below] at (1.5, 3) {$T_2$};
          \node [below] at (5.5, 1.7) {$T_1$};
          

\end{tikzpicture}
\end{document}
